\section{Finite repeated sampling on actual intervals}

The operation below uses a fresh coordinate at every active round and a
terminal full-family image and trim. It follows the resampling description
in \cite{Arad2016RigorousRG}; its distinction from the printed sample-once
procedure and the corrected constants are recorded in
\cite{gap:alvv16_bootstrap_constants}.

\begin{definition}[The complete finite-resampling interval parent]
    \label{def:complete_interval_resampling_parent}
    \lean{TNLean.MPS.AreaLaw.canonicalHalfSampledIntervalParent}\leanok
    Fix an interval $J=[p,r)$, operators $F_1,\ldots,F_n$ on
    $\mathcal H_J$, a threshold $\theta$, a trimming radius $\xi$, a
    parameter $D\ge2$ and a number of rounds $R$. Fix a coordinate isometry $I_J:\C^{d_J}\to\mathcal H_J$.
    Given a candidate $W$, start from its
    actual orthogonal projection $Q_0$ in those coordinates. At an active
    round $j$, where $\operatorname{rank}Q_j\ge\theta$, sample
    $\lfloor\operatorname{rank}Q_j/(2D^2)\rfloor$ dimensions of its actual
    range using the prescribed unitary coordinate for its actual rank. Apply the family
    to that sampled space and collectively trim at the actual internal
    cuts of $J$; let $Q_{j+1}$ be the resulting projection in the same
    finite coordinates. At an inactive round set $Q_{j+1}=Q_j$.

    After the $R$ prescribed coordinates, extend the sequence by identities.
    Transport the range of $Q_R$ to $\mathcal H_J$ and take its full family
    image $V_R=\sum_i F_i I_J\operatorname{ran}Q_R$. The complete parent is
    the collective trim of an orthonormal basis of this same $V_R$.
    In particular the terminal full image and trim are performed even
    when the initial candidate is already below the stopping threshold.
\end{definition}

\begin{lemma}[Measurability of the parent with a finite sequence]
    \label{lem:complete_interval_resampling_parent_measurable}
    \lean{TNLean.MPS.AreaLaw.measurable_projectionMatrix_canonicalHalfSampledIntervalParent}\leanok
    For every measurable candidate family with measurable orthogonal
    projection, the actual complete parent projection is jointly
    measurable in the candidate parameter and its finite seed sequence.
    No sampling-success hypothesis is required.
\end{lemma}
\begin{proof}\leanok
    \uses{def:complete_interval_resampling_parent}
    Reindex the actual candidate projection into the fixed finite
    coordinates. This is a continuous matrix reindexing. Extension of the
    finite sequence by identities is measurable, so composition with the
    joint measurable native parent projection proves the assertion.
\end{proof}

\begin{definition}[Independent finite sequences on the balanced tree]
    \label{def:interval_resampling_choices}
    \lean{TNLean.MPS.AreaLaw.IntervalResamplingChoices,
        TNLean.MPS.AreaLaw.intervalResamplingMeasure,
        TNLean.MPS.AreaLaw.instIsProbabilityMeasureIntervalResamplingMeasure,
        TNLean.MPS.AreaLaw.canonicalIntervalResamplingOperation}\leanok
    For each possible endpoint pair $I$, let $d_I=\dim\mathcal H_I$ and
    set $\Omega_I=(\prod_{a=0}^{d_I}\mathrm U(a))^R$, with its finite
    product Haar probability law $\mu_I$. Put
    $\Omega=\prod_I\Omega_I$ and $\mu=\bigotimes_I\mu_I$.
    Fix all interval families, thresholds and scalar parameters before
    $\omega\in\Omega$. At a valid long parent $J$, apply
    Definition~\ref{def:complete_interval_resampling_parent} to the actual
    adjacent child merge and the coordinate $\omega_J$.
    Full short leaves and the existing balanced recursion determine
    $S_\omega(p,r)$. No coordinate of an invalid endpoint pair is used.
\end{definition}

\begin{lemma}[Measurability and joint failure for the actual resampling tree]
    \label{lem:actual_interval_resampling_tree_probability}
    \lean{TNLean.MPS.AreaLaw.measurable_projectionMatrix_and_dependsOn_canonicalIntervalResamplingTree,
        TNLean.MPS.AreaLaw.measurableSet_isViable_canonicalIntervalResamplingTree,
        TNLean.MPS.AreaLaw.intervalResamplingMeasure_canonical_tree_joint_failure_le}\leanok
    Every valid $S_\omega(p,r)$ has measurable orthogonal projection and
    depends only on coordinates of no longer intervals. Its true viability
    event for a fixed global target and the whole exterior is measurable.

    Fix $\delta\ge0$. Suppose that for each long parent and every fixed
    baseline seed, viability of both actual children implies that the
    literal complete parent's failure has $\mu_J$-measure at most $\beta$.
    Then simultaneous failure of $\delta$-viability at a valid root or one
    of its actual listed internal intervals has outer measure at most
    $\max\{r-p-1,0\}\beta$.
\end{lemma}
\begin{proof}\leanok
    \uses{def:interval_resampling_choices,
        lem:complete_interval_resampling_parent_measurable,
        lem:independent_interval_tree_sections,
        lem:independent_interval_parent_probability}
    The actual complete parent satisfies the measurable-composition
    hypothesis by
    Lemma~\ref{lem:complete_interval_resampling_parent_measurable}.
    Apply the independent-coordinate projection and section estimates to
    this same parent and its finite sequence coordinate. Full short
    leaves are viable, and the guarded failure argument uses only the
    actual root merge list.
\end{proof}

\begin{lemma}[Unconditional recorded-rank bounds on the same tree]
    \label{lem:actual_interval_resampling_tree_profile}
    \lean{TNLean.MPS.AreaLaw.isIntervalSpanWithRankBound_canonicalHalfSampledIntervalParent,
        TNLean.MPS.AreaLaw.canonicalIntervalResamplingTree_profile}\leanok
    Fix $\xi>0$. For one parent, suppose $n\le D^2$, $\theta\ge2$ and
    $d_J\le2^R$. On every finite sequence, its output has at most
    $n\theta$ recorded generators, each having internal-cut Schmidt rank
    at most $n\theta/\xi^2$. Any larger count has the same corresponding
    bounds.

    For the balanced tree, suppose $q\ge1$, $0<\xi\le1$, the leaf width
    is positive, and a common $C$ satisfies $q^w\le C$ and
    $n_J\theta_J\le C$ at every valid long interval. Suppose also
    $n_J\le D^2$, $\theta_J\ge2$ and $d_J\le2^R$ there. Then every actual
    retained interval space has dimension at most $C$ and at most $C$
    recorded generators with weighted internal-cut ranks at most $C$,
    for every whole seed sequence.
\end{lemma}
\begin{proof}\leanok
    \uses{def:complete_interval_resampling_parent,def:interval_resampling_choices}
    The initial projection rank is at most $d_J\le2^R$. The stopped
    halving bound gives $\operatorname{rank}Q_R<\theta_J$ on every
    sequence. Thus the same full-family image has dimension at most
    $n_J\theta_J$. Identify that range with the actual sum of native
    family images and trim its standard orthonormal basis. The collective
    trimming estimate bounds both the recorded generator count and their
    weighted ranks by this image dimension. For the tree, induct on the
    balanced subdivision: a full short leaf fits the $q^w$ bound, and
    every actual parent satisfies the preceding unconditional estimate.
\end{proof}
